Fix an observed law \(p_{\mathrm w}\) satisfying the hypothesis, and let
\[
\mathcal F(p_{\mathrm w})
=\{M\in\mathfrak M_{\mathrm w,X}:p_M=p_{\mathrm w}\}
\]
be its model fiber.  We first prove the outer inclusion.  For
\(M\in\mathcal F(p_{\mathrm w})\), qualitative functionality gives a map
\(f:\{0,1\}\to\{0,1\}\) for which
\(\theta_A(a\mid z)=\mathbf 1\{a=f(z)\}\).  Put
\[
\pi_z=\theta_Z(z),\qquad
e_{ax}=\theta_X(x\mid a),\qquad
m_{ax}=\theta_Y(1\mid a,x).
\]
The recursive factorization assumption, applied to the edges of
\(G_{\mathrm w}\), and summation over the hidden value \(A\) give, for every
\(z,x\in\{0,1\}\),
\[
p_z=\pi_z,
\qquad
p_{zx}=\pi_z e_{f(z),x},
\qquad
p_{\mathrm w}(z,x,Y=1)=\pi_z e_{f(z),x}m_{f(z),x}.
\tag{1}
\]
By the truncated-intervention assumption, intervention on \(X\) deletes the
factor \(e_{f(z),x}\), and hence
\[
Q_{X,M}(x,1)=\sum_{z\in\{0,1\}}\pi_zm_{f(z),x}
=\sum_{z\in\{0,1\}}p_zm_{f(z),x}.
\tag{2}
\]
If \(p_{zx}>0\), all factors in the middle identity in (1) are nonnegative,
so \(p_z>0\) and \(e_{f(z),x}>0\); division in the last identity in (1) is
therefore licensed and yields
\[
m_{f(z),x}=\frac{p_{\mathrm w}(z,x,Y=1)}{p_{zx}}.
\tag{3}
\]
If \(p_z>0\) and \(p_{zx}=0\), then (1) instead implies
\(e_{f(z),x}=0\), while the Bernoulli row parameter only satisfies
\(0\le m_{f(z),x}\le1\).  Values with \(p_z=0\) make zero contribution to
(2).  Thus, with
\[
H_x=\{z:p_z>0,\ p_{zx}=0\},
\]
equations (2)--(3) and Definition~\texttt{def:binary-identified-region}
give
\[
Q_{X,M}(x,1)
=L_x+\sum_{z\in H_x}p_zm_{f(z),x},
\qquad
L_x\le Q_{X,M}(x,1)\le L_x+\sum_{z\in H_x}p_z
=L_x+M_x.
\tag{4}
\]
This argument remains valid when two values of \(z\) share the same latent
state: if \(f(z)=f(z')\) and both \(p_z,p_{z'}\) are positive, then (1) shows
that their \(X\)-support indicators agree because they use the same
\(e_{f(z),x}\).  Shared rows can restrict a particular model but cannot
invalidate (4).

We next prove simultaneous attainability and thereby the reverse inclusion.
Fix any
\[
(r_0,r_1)\in[L_0,L_0+M_0]\times[L_1,L_1+M_1],
\qquad
\delta_x=r_x-L_x\in[0,M_x].
\]
Construct a model \(M^*\) on the original four binary variables as follows.
Set
\[
f^*(z)=z,
\qquad
\theta_Z^*(z)=p_z,
\qquad
\theta_A^*(a\mid z)=\mathbf 1\{a=z\}.
\tag{5}
\]
For \(a=z\), define the two rows of \(X\) by
\[
\theta_X^*(x\mid a=z)=
\begin{cases}
p_{zx}/p_z,&p_z>0,\\
1/2,&p_z=0.
\end{cases}
\tag{6}
\]
The division in (6) is guarded by \(p_z>0\); moreover
\(\sum_xp_{zx}=p_z\), so each displayed row is normalized and nonnegative.
For the \(Y\)-rows put
\[
s_{zx}=
\begin{cases}
p_{\mathrm w}(z,x,Y=1)/p_{zx},&p_{zx}>0,\\
\delta_x/M_x,&z\in H_x\text{ and }M_x>0,\\
0,&\text{otherwise},
\end{cases}
\qquad
\theta_Y^*(1\mid a=z,x)=s_{zx},
\quad
\theta_Y^*(0\mid a=z,x)=1-s_{zx}.
\tag{7}
\]
Every division in (7) is explicitly guarded.  In its second case,
\(0\le\delta_x/M_x\le1\), so (7) defines Bernoulli rows.

We verify observational equality cell by cell.  If \(p_z=0\), both
\(p_{\mathrm w}(z,x,y)\) and the law induced by (5)--(7) are zero.  If
\(p_z>0\) and \(p_{zx}=0\), equation (6) makes both induced cells with that
\((z,x)\) zero.  Finally, if \(p_{zx}>0\), equations (5)--(7) give
\[
p_z\frac{p_{zx}}{p_z}
\begin{cases}s_{zx},&y=1,\\1-s_{zx},&y=0,
\end{cases}
=p_{\mathrm w}(z,x,y).
\tag{8}
\]
Thus \(p_{M^*}=p_{\mathrm w}\).  The node \(A\) is functional by (5), and
all mechanisms obey the factorization and truncated-intervention semantics
by construction.  Furthermore,
\[
P_{M^*}(X=x)=\sum_zp_{zx}=p_{\mathrm w}(X=x)>0
\]
for both \(x\), so the joint-treatment-positivity assumption is preserved and
\(M^*\in\mathfrak M_{\mathrm w,X}\).

Applying (2) to \(M^*\) and using (7), for each \(x\) with \(M_x>0\) we get
\[
Q_{X,M^*}(x,1)
=L_x+\sum_{z\in H_x}p_z\frac{\delta_x}{M_x}
=L_x+\delta_x=r_x.
\tag{9}
\]
If \(M_x=0\), then \(r_x=L_x\) and the same conclusion follows from (4)
with an empty contribution.  The rows indexed by \(x=0\) and \(x=1\) in
(7) are distinct, so the two choices are simultaneous.  Equations (4) and
(9) prove that the identified set is exactly the asserted Cartesian
rectangle and that every point is achieved without alphabet enlargement.

Finally, this rectangle is a singleton exactly when \(M_0=M_1=0\).  Because
the summands in \(M_x=\sum_{z\in H_x}p_z\) are strictly positive on \(H_x\),
\(M_x=0\) holds exactly when \(p_{zx}>0\) for every \(z\) with \(p_z>0\).
Applying this equivalence for both treatment values proves the stated
point-identification criterion.